\DAsection[镜像、容器、网络、卷与 Compose]{Docker 基础与部署模型}

\begin{frame}{7.1 Image 与 Container}{镜像是静态模板，容器是运行实例}
  \centering
  \begin{tikzpicture}[node distance=12mm]
    \node[DAflow,minimum width=34mm] (dockerfile) {Dockerfile\\构建说明};
    \node[DAflow,minimum width=34mm,right=of dockerfile] (image) {Image\\只读层 + 元数据};
    \node[DAflow,minimum width=34mm,right=of image] (container1) {Container A\\运行进程};
    \node[DAflow,minimum width=34mm,below=6mm of container1] (container2) {Container B\\另一个实例};
    \draw[DAarrowred] (dockerfile) -- node[above,font=\tiny]{docker build} (image);
    \draw[DAarrowred] (image) -- node[above,font=\tiny]{docker run} (container1);
    \draw[DAarrowred] (image) -- (container2);
  \end{tikzpicture}
  \par\medskip
  \begin{itemize}
    \item 容器的主进程退出，容器就停止。
    \item 容器内临时写入不会进入镜像，也不会自动持久化。
  \end{itemize}
  \DAsource{FastAPI Docs · What is a Container}{https://fastapi.tiangolo.com/deployment/docker/}
\end{frame}

\begin{frame}[fragile]{7.2 最小 Dockerfile 读法}{每条指令都在描述镜像构建过程}
  \begin{lstlisting}[language=bash]
FROM python:3.12-slim
WORKDIR /app
COPY requirements.txt .
RUN pip install --no-cache-dir -r requirements.txt
COPY . .
EXPOSE 8000
CMD ["python", "-m", "uvicorn", "app.main:app",
     "--host", "0.0.0.0", "--port", "8000"]
  \end{lstlisting}
  \begin{itemize}
    \item \DAcmd{RUN} 在构建时执行；\DAcmd{CMD} 在容器启动时执行。
    \item 容器内服务监听 \DAfile{0.0.0.0}，才能从容器网络访问。
    \item JSON Exec Form 能更正确地接收停止信号。
  \end{itemize}
\end{frame}

\begin{frame}{7.3 Docker 层与缓存}{把变化少的步骤放前面}
  \begin{enumerate}
    \item 复制依赖清单与锁文件。
    \item 安装依赖，生成可缓存层。
    \item 再复制变化频繁的业务源码。
    \item 执行构建，最后只把运行产物带进最终镜像。
  \end{enumerate}
  \begin{columns}[T]
    \begin{column}{0.48\textwidth}
      \DAstatbox{快}{依赖不变时复用缓存}
    \end{column}
    \begin{column}{0.48\textwidth}
      \DAstatbox{小}{多阶段构建不携带开发工具}
    \end{column}
  \end{columns}
  \DAtagline{常见错误}{先 \DAcmd{COPY . .} 再安装依赖，任何源码变化都会让依赖层缓存失效}
\end{frame}

\begin{frame}[fragile]{7.4 .dockerignore 控制构建上下文}{缩小构建上下文并排除敏感文件}
  \begin{lstlisting}
.git
.env
node_modules
.pnpm-store
dist
.next
.venv
__pycache__
*.log
coverage
  \end{lstlisting}
  \begin{itemize}
    \item Docker 构建只应该接收生成镜像所需的文件。
    \item 即使 Dockerfile 没有 COPY 某文件，发送巨大构建上下文仍会拖慢构建。
    \item \DAfile{.env}、私钥和本地虚拟环境必须排除。
  \end{itemize}
\end{frame}

\begin{frame}[fragile]{7.5 Build、Run、Port 三个动作}{左边是宿主机端口，右边是容器端口}
  \begin{lstlisting}[language=bash]
docker build -t customer-api:1.0 .

docker run -d \
  --name customer-api \
  -p 127.0.0.1:8000:8000 \
  --restart unless-stopped \
  customer-api:1.0
  \end{lstlisting}
  \begin{itemize}
    \item \DAfile{127.0.0.1:8000:8000} 只绑定宿主机本地，供 Nginx 反代。
    \item 如果写 \DAfile{8000:8000}，通常会绑定全部网卡，应再由防火墙限制。
    \item Image Tag 是版本坐标；生产环境使用明确版本或 Commit SHA。
  \end{itemize}
\end{frame}

\begin{frame}[fragile]{7.6 容器排错的四个命令}{状态、日志、进程内检查、配置}
  \begin{lstlisting}[language=bash]
docker ps -a
docker logs --tail 200 -f customer-api
docker exec -it customer-api sh
docker inspect customer-api
  \end{lstlisting}
  \begin{itemize}
    \item \DAcmd{docker ps -a} 看退出状态和重启循环。
    \item \DAcmd{docker logs} 看应用标准输出与错误输出。
    \item \DAcmd{docker exec} 用于诊断，不应在容器内手工修改生产代码。
  \end{itemize}
\end{frame}

\begin{frame}{7.7 Volume 解决数据持久化}{容器可替换，数据必须保留}
  \begin{DAtable}{@{}p{31mm}p{40mm}p{51mm}@{}}
    \toprule
    数据 & 是否放镜像 & 推荐位置 \\
    \midrule
    应用源码 / 构建物 & 是 & Image 只读层 \\
    MySQL 数据目录 & 否 & Named Volume / 云数据库 \\
    用户上传文件 & 否 & Volume / 对象存储 \\
    日志 & 通常否 & 标准输出 + 日志平台 / 受控目录 \\
    Secret & 否 & 环境变量 / Secret 文件 \\
    \bottomrule
  \end{DAtable}
  \DAtagline{设计原则}{应用容器应尽量无状态，方便重建、替换和横向扩展}
\end{frame}

\begin{frame}{7.8 Docker Network 让服务用名字互访}{Compose 会为项目创建默认网络}
  \centering
  \begin{tikzpicture}[node distance=9mm]
    \node[DAflow,minimum width=31mm] (frontend) {frontend\\Nginx};
    \node[DAflow,minimum width=31mm,right=of frontend] (api) {api\\FastAPI:8000};
    \node[DAstore,right=of api] (db) {db\\MySQL:3306};
    \draw[DAarrowred] (frontend) -- node[above,font=\tiny]{http://api:8000} (api);
    \draw[DAarrow] (api) -- node[above,font=\tiny]{mysql://db:3306} (db);
  \end{tikzpicture}
  \par\medskip
  \begin{itemize}
    \item 容器间不使用 \DAfile{localhost}；\DAfile{localhost} 只代表当前容器自己。
    \item 使用服务名 \DAfile{api}、\DAfile{db} 作为网络主机名。
    \item 只有需要给宿主机 / 公网访问的服务才映射 ports。
  \end{itemize}
\end{frame}

\begin{frame}{7.9 Healthcheck 与 restart}{分别负责健康判断与重启策略}
  \begin{DAtable}{@{}p{29mm}p{43mm}p{50mm}@{}}
    \toprule
    机制 & 判断对象 & 用途 \\
    \midrule
    restart policy & 主进程是否退出 & 崩溃或重启服务器后恢复容器 \\
    healthcheck & 应用是否能响应 & 标记 unhealthy、辅助依赖与监控 \\
    readiness & 是否可接流量 & 大型编排平台控制流量进入 \\
    \bottomrule
  \end{DAtable}
  \begin{itemize}
    \item 健康检查应快速、稳定，不执行昂贵业务。
    \item 可检查应用事件循环、关键依赖连接，但避免产生副作用。
  \end{itemize}
\end{frame}

\begin{frame}[fragile]{7.10 Docker Compose 管理多容器应用}{一份 YAML 描述服务、网络、卷与配置}
  \begin{lstlisting}[language=bash]
docker compose config
docker compose build
docker compose up -d
docker compose ps
docker compose logs -f
docker compose down
  \end{lstlisting}
  \begin{itemize}
    \item \DAcmd{config} 先检查合并后的配置和变量替换。
    \item \DAcmd{down} 默认不删 Named Volume；加 \DAcmd{-v} 会删除数据卷，必须谨慎。
    \item Compose 适合单机部署，也是今天全栈 Docker 路线的主角。
  \end{itemize}
  \DAsource{Docker Docs · Compose}{https://docs.docker.com/compose/}
\end{frame}
